\documentclass[11pt,a4paper]{article}
\usepackage[utf8]{inputenc}
\usepackage[spanish]{babel}
\usepackage{graphicx}
\usepackage{geometry}
\usepackage{hyperref}
\usepackage{xcolor}
\usepackage{fancyhdr}
\usepackage{enumitem}
\usepackage{booktabs}
\usepackage{longtable}
\usepackage{array}
\usepackage{colortbl}
\usepackage{tabularx}
\usepackage{mdframed}
\usepackage{float}

\definecolor{azuloscuro}{RGB}{44,82,130}
\definecolor{azulclaro}{RGB}{235,244,255}
\definecolor{grisclaro}{RGB}{247,247,247}
\definecolor{verdeclaro}{RGB}{240,255,244}
\definecolor{rojoclaro}{RGB}{255,245,245}
\definecolor{azulinfo}{RGB}{235,248,255}

\geometry{top=2.5cm,bottom=2.5cm,left=3cm,right=3cm}
\pagestyle{fancy}
\fancyhf{}
\fancyhead[L]{\small IF614 -- Ingeniería de Software I}
\fancyhead[R]{\small Residuos Cusco -- Entrega 2}
\fancyfoot[C]{\thepage}

\hypersetup{colorlinks=true, linkcolor=azuloscuro, urlcolor=azuloscuro}

\newcommand{\theadcell}[1]{%
  \cellcolor{azuloscuro}\textcolor{white}{\textbf{\small #1}}}

\begin{document}

% =================================================================
% CARÁTULA
% =================================================================
\thispagestyle{empty}
\begin{center}
    \Large \textbf{UNIVERSIDAD NACIONAL DE SAN ANTONIO ABAD DEL CUSCO}\\[0.5cm]
    \large \textbf{FACULTAD DE INGENIERÍA ELÉCTRICA, ELECTRÓNICA, INFORMÁTICA Y MECÁNICA}\\[0.5cm]
    \large \textbf{ESCUELA PROFESIONAL DE INGENIERÍA INFORMÁTICA Y DE SISTEMAS}\\[0.3cm]
    \noindent\rule{\textwidth}{1.5pt}\\[0.5cm]
    \begin{minipage}{0.45\textwidth}
        \begin{center}
            \includegraphics[width=0.7\textwidth]{LOGOUNSAAC.png}
        \end{center}
    \end{minipage}
    \hfill
    \begin{minipage}{0.45\textwidth}
        \begin{center}
            \includegraphics[width=0.7\textwidth]{LOGOCARRERA.png}
        \end{center}
    \end{minipage}\\[0.3cm]
    \noindent\rule{\textwidth}{0.5pt}\\[0.2cm]
    \Huge \textbf{SISTEMA INTELIGENTE DE RECOLECCIÓN}\\[0.1cm]
    \Huge \textbf{DE RESIDUOS SÓLIDOS SEGREGADOS}\\[0.2cm]
    \Large \textbf{Entrega N.° 2 -- Documentación del Proyecto}\\[0.2cm]
    \large \textbf{Curso:} IF614 -- Ingeniería de Software I\\[0.3cm]
    \large \textbf{Docente(s):} Gamarra Salas, Jisbaj; Cosio Loaiza, Stephan Jhoel; Ticona Felix, Luz Indira\\[0.2cm]
    \large \textbf{Integrantes:}\\[0.4cm]
    \renewcommand{\arraystretch}{1}
    \begin{tabular}{|p{7cm}|p{2.5cm}|p{3.5cm}|}
    \hline
    \rowcolor{azuloscuro}
    \textcolor{white}{\textbf{Apellidos y Nombres}} &
    \textcolor{white}{\textbf{Código}} &
    \textcolor{white}{\textbf{Rol}} \\
    \hline
    Luicho Quispe, Jhoel Alex        & 224871 & Product Owner \\[0.2cm]
    \hline
    Luna Ccapa, Carlos Willian       & 210178 & Developer \\[0.2cm]
    \hline
    Puma Condori, Richard Braulio    & 161809 & Developer \\[0.2cm]
    \hline
    Huaracallo Arenas, Lino Zeynt    & 204798 & Scrum Master \\[0.2cm]
    \hline
    \end{tabular}\\[0.5cm]
    \large Metodología: SCRUM \quad $|$ \quad Atributo del graduado: AG-C01\\[0.3cm]
    \large CUSCO -- PERÚ\\[0.2cm]
    \large 2026
\end{center}
\newpage

\tableofcontents
\newpage

% =================================================================
% CAPÍTULO I
% =================================================================
\section{Datos Generales del Proyecto}

\subsection{Título del proyecto}
Diseño e Implementación de un Sistema Inteligente de Recolección de Residuos Sólidos Segregados para la Gestión Ambiental Urbana en la ciudad del Cusco.

\subsection{Objetivo general y específicos}
\textbf{Objetivo general:} diseñar e implementar un sistema inteligente de gestión de residuos sólidos segregados que optimice la recolección, mejore la comunicación con los ciudadanos y genere información estratégica para la toma de decisiones, aplicando la metodología ágil SCRUM.

\textbf{Objetivos específicos:}
\begin{enumerate}[leftmargin=*]
    \item Desarrollar el módulo de gestión de usuarios y zonas con registro de ciudadanos y asignación geográfica (detección por GPS).
    \item Implementar un catálogo de tipos de residuos con guía de segregación (orgánico, reciclable, no reciclable, peligroso).
    \item Construir el monitoreo de rutas con seguimiento de la ubicación de los camiones.
    \item Ofrecer una aplicación móvil para consultar rutas y reportar incidencias con geolocalización.
    \item Proveer un sistema de alertas y un módulo de reportes para la municipalidad.
\end{enumerate}

\subsection{Alcance del sistema}
El sistema permite a la municipalidad y a los ciudadanos del Cusco gestionar el proceso de recolección de residuos sólidos segregados: administración de usuarios y zonas, catálogo de residuos, monitoreo de rutas, aplicación móvil ciudadana, sistema de alertas y módulo de reportes. \textit{Fuera de alcance:} hardware IoT en contenedores, pasarelas de pago e integración con sistemas tributarios.

\subsection{Justificación del proyecto}
La gestión de residuos en el Cusco presenta rutas no optimizadas, escasa comunicación con la población y baja cultura de segregación, lo que genera acumulación de residuos y afecta la salubridad y la sostenibilidad ambiental. Una solución digital con monitoreo, comunicación ciudadana y analítica contribuye al desarrollo sostenible y a una gestión municipal más eficiente y transparente.

\subsection{Tecnologías utilizadas}
\begin{table}[H]
\centering
\renewcommand{\arraystretch}{1.3}
\begin{tabularx}{\textwidth}{|l|X|}
\hline
\theadcell{Capa} & \theadcell{Tecnología} \\
\hline
Backend / API & Node.js 20 + Express \\
\hline
Base de datos & MongoDB (MongoDB Atlas) + Mongoose \\
\hline
Autenticación & JWT (jsonwebtoken) + bcryptjs \\
\hline
Aplicación móvil & Expo SDK 54 + React Native + expo-router (TypeScript) \\
\hline
Dashboard web & HTML + JavaScript (servido por el backend) \\
\hline
Mapas / GPS & OpenStreetMap (Leaflet) + expo-location \\
\hline
Integraciones & API RENIEC (apisperu) para autocompletar por DNI \\
\hline
Despliegue & Render (backend + web) · APK (Android) \\
\hline
Gestión / colaboración & Git · GitHub \\
\hline
\end{tabularx}
\end{table}

\subsection{Roles del equipo Scrum}
\begin{table}[H]
\centering
\renewcommand{\arraystretch}{1.3}
\begin{tabularx}{\textwidth}{|l|l|X|}
\hline
\theadcell{Rol} & \theadcell{Integrante} & \theadcell{Funciones} \\
\hline
Scrum Master & Huaracallo Arenas, Lino Zeynt (204798) & Facilita el proceso, remueve impedimentos, asegura las ceremonias Scrum. \\
\hline
Product Owner & Luicho Quispe, Jhoel Alex (224871) & Prioriza el Product Backlog, define criterios de aceptación, representa al cliente. \\
\hline
Developer (Full Stack) & Luna Ccapa, Carlos Willian (210178) & Backend, app móvil y dashboard; pruebas. \\
\hline
Developer (Full Stack) & Puma Condori, Richard Braulio (161809) & Backend, app móvil y dashboard; pruebas. \\
\hline
Developer (Full Stack) & Luicho Quispe, Jhoel Alex (224871) & Integración y despliegue. \\
\hline
\end{tabularx}
\end{table}
\noindent\textit{Nota:} todos los integrantes participan como desarrolladores full stack; Lino Huaracallo asume además el rol de Scrum Master y Jhoel Luicho el de Product Owner.

\subsection{Público objetivo del sistema}
\begin{itemize}[leftmargin=*]
    \item \textbf{Ciudadanos:} consultan rutas y horarios, reciben alertas, reportan incidencias y aprenden a segregar.
    \item \textbf{Operadores municipales:} ejecutan rutas y reportan su ubicación.
    \item \textbf{Administradores municipales:} gestionan zonas, rutas y usuarios, y consultan reportes.
\end{itemize}

% =================================================================
% CAPÍTULO II
% =================================================================
\newpage
\section{Método SCRUM Aplicado}

\subsection{Product Backlog}
\textbf{Requisitos Funcionales (priorizados):}
\begin{table}[H]
\centering
\renewcommand{\arraystretch}{1.25}
\begin{tabularx}{\textwidth}{|l|X|c|c|}
\hline
\theadcell{ID} & \theadcell{Requisito funcional} & \theadcell{Prioridad} & \theadcell{Estado} \\
\hline
RF-01 & Registro de ciudadanos & Alta & Completo \\
\hline
RF-02 & Autenticación e inicio de sesión (JWT) & Alta & Completo \\
\hline
RF-03 & Gestión de zonas de recolección & Alta & Completo \\
\hline
RF-04 & Asignación de usuarios a zonas (detección GPS) & Media & Completo \\
\hline
RF-05 & Registro de tipos de residuos & Media & Completo \\
\hline
RF-06 & Clasificación de residuos por categoría & Alta & Completo \\
\hline
RF-07 & Visualización de rutas de recolección & Alta & Completo \\
\hline
RF-08 & Seguimiento de ubicación de camiones & Alta & Completo \\
\hline
RF-09 & Gestión de rutas por el administrador & Alta & Completo \\
\hline
RF-10 & Consulta de rutas/horarios (app móvil) & Alta & Completo \\
\hline
RF-11 & Reporte ciudadano de incidencias (GPS) & Alta & Completo \\
\hline
RF-12 & Notificación de cercanía del camión & Alta & Parcial \\
\hline
RF-13 & Alertas de retraso o incidencias & Media & Parcial \\
\hline
RF-14 & Reporte de residuos recolectados por zona & Alta & Completo \\
\hline
RF-15 & Reporte de cumplimiento de rutas & Media & Parcial \\
\hline
RF-16 & Reporte de participación ciudadana & Media & Parcial \\
\hline
\end{tabularx}
\end{table}

\textbf{Requisitos No Funcionales:}
\begin{itemize}[leftmargin=*]
    \item \textbf{RNF-01 Seguridad:} contraseñas cifradas (bcrypt) y autenticación por token JWT.
    \item \textbf{RNF-02 Disponibilidad:} backend y base de datos desplegados en la nube (24/7).
    \item \textbf{RNF-03 Usabilidad:} interfaz simple y accesible.
    \item \textbf{RNF-04 Rendimiento:} respuestas comunes en menos de 3 s.
    \item \textbf{RNF-05 Portabilidad:} app instalable en Android (APK).
    \item \textbf{RNF-06 Privacidad:} cumplimiento de la Ley N.° 29733 de Protección de Datos Personales.
    \item \textbf{RNF-07 Mantenibilidad:} código modular y versionado en GitHub.
\end{itemize}

\subsection{Sprints realizados}
Número total de sprints: \textbf{3} · Duración: \textbf{2 semanas} cada uno.
\begin{table}[H]
\centering
\renewcommand{\arraystretch}{1.3}
\begin{tabularx}{\textwidth}{|c|X|X|}
\hline
\theadcell{Sprint} & \theadcell{Objetivo} & \theadcell{Entregables} \\
\hline
1 & Autenticación y gestión de usuarios/zonas & API de auth (JWT), modelos, zonas, detección por GPS \\
\hline
2 & Núcleo operativo y app móvil & Rutas + ubicación, residuos, incidencias, pantallas móviles \\
\hline
3 & Alertas, reportes, web y despliegue & Dashboard web, reportes, despliegue en la nube, APK \\
\hline
\end{tabularx}
\end{table}

\subsection{Herramientas de colaboración}
GitHub (repositorio y ramas por historia de usuario), Render (despliegue del backend y la web), MongoDB Atlas (base de datos en la nube) y Expo/EAS (desarrollo y compilación de la app móvil).

% =================================================================
% CAPÍTULO III (BOSQUEJO)
% =================================================================
\newpage
\section{Desarrollo por Sprint (Bosquejo)}

\subsection{Sprint 1 — Autenticación y gestión de usuarios/zonas}
\begin{itemize}[leftmargin=*]
    \item \textbf{Objetivo:} habilitar registro/login seguro y la gestión de zonas con asignación por ubicación.
    \item \textbf{Historias:} RF-01, RF-02, RF-03, RF-04 (completadas).
    \item \textbf{Tareas:} modelos Usuario/Zona, hash de contraseñas, emisión de JWT, endpoint de detección de zona (punto-en-polígono).
    \item \textbf{Entregables:} API de autenticación funcional y datos de prueba (seed).
    \item \textbf{Diagramas:} casos de uso (Módulo 1), clases (Usuario, Zona).
    \item \textbf{Sprint Review:} validación de login/registro; \textit{feedback:} agregar autocompletado por DNI.
    \item \textbf{Retrospective:} reforzar pruebas del endpoint de zonas.
\end{itemize}

\subsection{Sprint 2 — Núcleo operativo y app móvil}
\begin{itemize}[leftmargin=*]
    \item \textbf{Objetivo:} implementar rutas, residuos e incidencias, y las primeras pantallas móviles.
    \item \textbf{Historias:} RF-05 a RF-11 (completadas).
    \item \textbf{Entregables:} CRUD de rutas + ubicación, catálogo de residuos, incidencias con GPS, app móvil navegable.
    \item \textbf{Diagramas:} casos de uso (Módulos 2-4), secuencia (reporte de incidencia), actividades (registro).
    \item \textbf{Sprint Review:} validación del reporte de incidencias; \textit{feedback:} detección de zona en el registro.
    \item \textbf{Retrospective:} unificar el formato de respuestas de la API.
\end{itemize}

\subsection{Sprint 3 — Alertas, reportes, web y despliegue}
\begin{itemize}[leftmargin=*]
    \item \textbf{Objetivo:} cerrar funcionalidades y desplegar el sistema.
    \item \textbf{Historias:} RF-12 a RF-16, dashboard web, despliegue y APK.
    \item \textbf{Entregables:} alertas, reportes por zona/categoría, dashboard web, backend en Render, APK Android.
    \item \textbf{Diagramas:} componentes (arquitectura), despliegue.
    \item \textbf{Sprint Review:} demostración del sistema completo (web + móvil).
    \item \textbf{Retrospective:} documentar y preparar la exposición.
\end{itemize}

% =================================================================
% ARQUITECTURA
% =================================================================
\newpage
\section{Diseño y Arquitectura del Sistema}
\subsection{Arquitectura general}
Arquitectura \textbf{cliente--servidor} en la nube:
\begin{itemize}[leftmargin=*]
    \item \textbf{Clientes:} app móvil (Android/APK) y dashboard web (navegador).
    \item \textbf{Servidor:} API REST (Node.js + Express) desplegada en Render.
    \item \textbf{Datos:} MongoDB Atlas.
    \item Comunicación HTTPS/JSON con autenticación JWT.
\end{itemize}
\begin{center}
\textit{App móvil / Web $\longrightarrow$ API REST (Render) $\longrightarrow$ MongoDB Atlas}
\end{center}

\subsection{Modelo de datos (colecciones MongoDB)}
Usuario, Zona (polígono geográfico), Ruta (estado y ubicación en vivo), Incidente, Residuo y Alerta. Relaciones: el Usuario pertenece a una Zona; la Ruta cubre una Zona y tiene un operador; el Incidente es reportado por un Usuario; el Residuo se registra por Zona.

\subsection{Patrones y mantenibilidad}
Separación por capas (rutas, modelos, middleware), respuestas uniformes \texttt{\{success, message, data\}}, control de acceso por rol y versionado en GitHub con ramas por historia. El sistema escala añadiendo zonas/rutas y desplegando réplicas del backend.

% =================================================================
% RESULTADOS Y CONCLUSIONES
% =================================================================
\newpage
\section{Resultados y Conclusiones}
\subsection{Resultados obtenidos}
\begin{itemize}[leftmargin=*]
    \item Sistema de tres capas funcional y \textbf{desplegado en la nube} (supera el 50\,\% exigido para la Entrega 2).
    \item \textbf{11 de 16} requisitos funcionales completos; los módulos esenciales operativos.
    \item App móvil instalable (APK) e integración con la API RENIEC (consulta de DNI).
\end{itemize}

\subsection{Impacto sostenible (AG-C01)}
\begin{itemize}[leftmargin=*]
    \item \textbf{Ambiental:} promueve la segregación y la recolección eficiente de residuos.
    \item \textbf{Social:} fomenta la participación ciudadana y la comunicación con la municipalidad.
    \item \textbf{Legal/ético:} cumplimiento de la Ley N.° 29733 (datos personales), Ley N.° 27314 (residuos sólidos) y NTP 900.058 (código de colores).
\end{itemize}

\subsection{Dificultades y soluciones}
\begin{itemize}[leftmargin=*]
    \item \textit{Pruebas en red local inestables} $\rightarrow$ se desplegó el backend en la nube (Render) con URL permanente.
    \item \textit{Detección de zona limitada} $\rightarrow$ se ampliaron los polígonos para cubrir el área metropolitana.
\end{itemize}

\subsection{Conclusiones}
\begin{enumerate}[leftmargin=*]
    \item Se obtuvo un sistema funcional (backend, app móvil y dashboard web) desplegado en la nube.
    \item Se aplicó SCRUM con Product Backlog priorizado, 3 sprints y ramas por historia de usuario.
    \item El proyecto atiende un problema real con impacto ambiental, social y cumplimiento legal.
    \item \textbf{Próximos pasos:} notificaciones push (RF-12/13), reportes avanzados (RF-15/16) y refinamiento del dashboard.
\end{enumerate}

\subsection*{Repositorio}
\url{https://github.com/GamoLG/Gestion-de-Residuos-Cusco}

\end{document}
